
% !TEX root = ../Testa/Principale.tex

\usepackage{iftex}

\ifInCompiler\relax\else
    \usepackage[
        pass,
        showframe=\bozza
    ]{geometry}\fi

\usepackage{amsmath}
\usepackage{amsthm}
\ifInCompiler\relax\else
    % \usepackage{amssymb}
    \usepackage{bbold}
\fi

\ifLuaTeX
    \usepackage{fontspec}
\else
    \usepackage[T1]{fontenc}
    % \usepackage[utf8]{inputenc}
    \usepackage[utf8]{inputenx}
\fi
\usepackage{microtype}

\usepackage[italian]{babel}

\usepackage{csquotes}
\usepackage[
    backend=biber,
    style=numeric-comp
]{biblatex}
\DeclareDelimFormat{finalnamedelim}{\addspace\&\space}
\addbibresource{../Corpo/Bibliografia.bib}
\nocite{*}

\usepackage[dvipsnames]{xcolor}

\ifInCompiler\relax\else
    \usepackage[
        drafting=\bozza,
        eulerchapternumbers,
        % subfig,
        beramono,
        eulermath,
        pdfspacing,
        dottedtoc=true
    ]{classicthesis}
    \usepackage{arsclassica}
    \usepackage{scrhack}
\fi
\usepackage[strict]{changepage}

\usepackage{tikz}
\usetikzlibrary{
    positioning,
    shapes.geometric,
    plotmarks,
    arrows.meta,
    calc,
    decorations.pathreplacing,
    backgrounds,
    patterns.meta,
    fpu, % «floating point unit»
    % external,
} 
\usepackage{shellesc}
\immediate\write18{if not exist "../../Figure/.fnl/TikZpdf" mkdir "../../Figure/.fnl/TikZpdf"}

\usepackage{etoolbox}
\newbool{forcetikzrecompile}

% \tikzexternalize[prefix={../../Figure/.fnl/TikZpdf/}]
% \tikzset{
%     % external/force remake,
%     external/system call={pdflatex \tikzexternalcheckshellescape -interaction=batchmode -jobname "\image" "\texsource"}
% }
% Il comando «\cref{}» non funziona all'interno di «\begin{tikzpicture}» se si attiva «\tikzexternalize»
% https://tex.stackexchange.com/questions/245411/using-ref-in-tikzpicture-node-while-using-externalize

\usepackage{pgfplots}
\pgfplotsset{compat=newest}
\usepgfplotslibrary{
    patchplots,
    groupplots,
    colorbrewer
}
\pgfdeclarelayer{background}
\pgfdeclarelayer{foreground}
\pgfsetlayers{background,main,foreground}

\usepackage{grffile}
\usepackage{nicematrix}

\usepackage{zref-savepos}
\newcounter{figposcount} % Crea un contatore unico per tracciare la posizione
\usepackage{xstring}
    
\usepackage{accents}

\usepackage{newunicodechar}

\usepackage{
    graphicx,
    caption,
    subcaption,
    float,
    stfloats
}

\usepackage[left]{lineno}

\usepackage[
    italiano,
    italianokw,
    ruled,
    vlined,
    linesnumbered
]{algorithm2e}
% \usepackage{listings}
\usepackage[
    section,
    newfloat
]{minted} % Va caricato dopo «lineno»

\usepackage{enumitem}
\usepackage{ragged2e} 

\ifInCompiler
    \usepackage{xr-hyper}
    \usepackage{hyperref} % È già caricato da «classicthesis»
\else
    \PassOptionsToPackage{hyphens}{url}
    \hypersetup{
        colorlinks=true,
        linktocpage=true,
        breaklinks=true,
        urlcolor=CTurl,
        linkcolor=CTlink,
        citecolor=CTcitation,
        % destlabel=true % Migliora la precisione dei nessi ipertestuali
    } % Si v. la linea 732 di «classicthesis.sty»
    % \hypersetup{
    %     colorlinks=false,
    %     linkbordercolor=CTlink,
    %     citebordercolor=CTcitation,
    %     urlbordercolor=CTurl
    %     linktocpage=true,
    %     breaklinks=true,
    %     pdfborder={1 1 1},
    % }
    % \usepackage[all]{hypcap}
\fi
\usepackage{xurl}

\usepackage[italian]{cleveref}

\usepackage{siunitx}
\sisetup{
    exponent-mode=threshold,
    exponent-thresholds=-3:3,
    round-mode=places,% figures
    round-precision=3,
    round-pad=false,
    separate-uncertainty=true,
    print-unity-mantissa=false,
    tight-spacing=true,
}

\usepackage[normalem]{ulem}
\usepackage{circledsteps}

\usepackage{afterpage}
\usepackage{booktabs}
\usepackage{needspace}
